%% 09-appendices.tex

\titleformat{\section}
  {\normalfont\normalsize\bfseries\scshape\color{navy}}
  {}{0em}{\Needspace*{4\baselineskip}Appendix~\thesection.\enspace}
  [\vspace{1.5pt}{\color{navy}\titlerule[0.4pt]}]

% ─────────────────────────────────────────────────────────────────────────────
\section{Contributions}
\label{app:contributions}

A breakdown of team contributions, encompassing both implementation and manuscript authorship, is provided in Table~\ref{tab:contributions}. At the request of collaborators at SickKids, the project's source code is hosted privately at \url{https://github.com/biswurmt/mlmd-noise/}. We plan to grant the instruction team access upon receipt of their GitHub handles.


\begin{table}[ht]
\centering
\caption{Detailed breakdown of teammate contributions to this report and underlying investigation.}
\label{tab:contributions}
\small
\begin{tabularx}{\textwidth}{@{} l X @{}}
\toprule
\textbf{Section} & \textbf{Contributors} \\
\midrule
Introduction          & Tyler: narrative structure and clinical framing \newline
                        Alice: technical enrichment and editing \\
\addlinespace
Related work          & Tyler: literature review and source curation across both tracks; synthesis and final write-up for noise-resistant training track \newline
                        Alice: literature review, source discovery, synthesis, and final write-up for label imputation track \\
\addlinespace
Data \& setting       & Alice: cohort selection, exclusions, and feature space description \\
\addlinespace
Methods               & \\
\addlinespace
\hspace{1em}Label imputation      & Alice: knowledge graph construction, multi-agent audit pipeline, and imputation strategies \newline
                                    Tyler: ideation and design (RAG strategy over EPFL dataset, bidirectional KG prediction, etc.) \\
\addlinespace
\hspace{1em}Noise-resistant training & Tyler: LiLAW implementation; proxy experimental design and validation; ablation analysis \\
\addlinespace
\hspace{1em}Diagnotix             & Alice: knowledge graph visualization, Navigate Mode, and dynamic pathway generation \newline
                                    Tyler: Analyze Mode chatbot, grounded via RAG over embedded clinical chunks \\
\addlinespace
Results               & \\
\addlinespace
\hspace{1em}KG validation \& discoveries & Alice: KG clinical validation and discovery of clinical false negatives \\
\addlinespace
\hspace{1em}MLMD performance      & Alice: impact of label imputation \newline
                                    Tyler: proxy validation and combined efficacy \\
\addlinespace
\midrule
Discussion            & Alice: primary synthesis and label imputation interpretation \newline
                        Tyler: noise-resistant training discussion and limitations \\
\addlinespace
Future work           & Alice \\
\bottomrule
\end{tabularx}
\end{table}

% ─────────────────────────────────────────────────────────────────────────────
\newpage
\section{Ontology Enrichment via External APIs}
\label{app:api_integration}

To ensure interoperability and standardize clinical entities, the Knowledge
Graph build pipeline queries eight external services to attach standard medical
codes and evidence metadata to each extracted record.
Table~\ref{tab:api_breakdown} outlines this sequential enrichment process; step
numbers correspond to the nine-step build pipeline described in the Methods
section.

\begin{table}[htbp]
    \centering
    \caption{API integration steps, targeted ontologies, and resulting dataset
    schema modifications.}
    \label{tab:api_breakdown}
    \small
    \resizebox{\linewidth}{!}{%
    \begin{tabular}{c p{0.26\linewidth} p{0.14\linewidth} p{0.20\linewidth} p{0.26\linewidth}}
        \toprule
        \textbf{Step} & \textbf{API (Target Ontology)} & \textbf{Credential} &
        \textbf{Enriched Entity} & \textbf{Column(s) Added} \\
        \midrule
        2 & \href{https://www.ebi.ac.uk/ols4}{EMBL-EBI OLS4} (HP / MONDO / EFO)
          & None
          & \texttt{raw\_symptom}
          & \texttt{ebi\_name}, \texttt{ebi\_code}, \texttt{synonyms} \\
        \addlinespace
        3 & \href{https://infocentral.infoway-inforoute.ca}{Canada Health Infoway FHIR}
            (SNOMED CT CA)
          & OAuth2 client credentials
          & \texttt{raw\_symptom}
          & \texttt{infoway\_name}, \texttt{infoway\_code} \\
        \addlinespace
        4 & \href{https://documentation.uts.nlm.nih.gov/rest/home.html}{UMLS REST API}
            (LOINC)
          & \texttt{UMLS\_API\_KEY}
          & \texttt{test} (unique values only)
          & \texttt{loinc\_code} \\
        \addlinespace
        5 & \href{https://documentation.uts.nlm.nih.gov/rest/home.html}{UMLS REST API}
            (ICD-10, WHO international)
          & \texttt{UMLS\_API\_KEY}
          & \texttt{condition} (unique values only)
          & \texttt{icd10\_code} \\
        \addlinespace
        6 & \href{https://infocentral.infoway-inforoute.ca}{Canada Health Infoway FHIR}
            (ICD-10-CA)
          & OAuth2 client credentials
          & \texttt{condition} (unique values only)
          & \texttt{icd10ca\_code} \\
        \addlinespace
        7 & \href{https://lhncbc.nlm.nih.gov/RxNav/APIs/index.html}{NLM RxNorm} \&
            \href{https://open.fda.gov/apis/}{OpenFDA}
          & None
          & \texttt{treatment} (where present)
          & \texttt{rxcui}, \texttt{adverse\_event} \\
        \addlinespace
        8 & \href{https://europepmc.org/RestfulWebService}{Europe PMC REST API}
          & None
          & \texttt{symptom} $\rightarrow$ \texttt{condition}
            (\texttt{INDICATES\_CONDITION})
          & \texttt{literature\_articles}, \texttt{literature\_patents} \\
        \addlinespace
        9 & \href{https://europepmc.org/RestfulWebService}{Europe PMC REST API}
          & None
          & \texttt{condition} $\rightarrow$ \texttt{test}
            (\texttt{REQUIRES\_TEST})
          & \texttt{literature\_articles}, \texttt{literature\_patents} \\
        \addlinespace
        10 & \href{https://europepmc.org/RestfulWebService}{Europe PMC REST API}
          & None
          & \texttt{symptom} $\rightarrow$ \texttt{test}
            (\texttt{DIRECTLY\_\allowbreak{}INDICATES\_\allowbreak{}TEST})
          & \texttt{literature\_articles}, \texttt{literature\_patents} \\
        \addlinespace
        11 & \href{https://clinicaltrials.gov/api/gui}{ClinicalTrials.gov v2}
          & None
          & \texttt{condition} $\rightarrow$ \texttt{treatment}
            (\texttt{RECOMMENDS\_TREATMENT})
          & \texttt{trial\_count} \\
        \bottomrule
    \end{tabular}}
\end{table}

% ─────────────────────────────────────────────────────────────────────────────
\newpage
\section{Multi-Agent Graph Audit Pipeline}
\label{app:audit}

\subsection*{Audit Unit and Triad Extraction}

The unit of audit is a clinical triad: a (finding $\rightarrow$ condition
$\rightarrow$ test) path extracted by traversing \texttt{INDICATES\_CONDITION}
followed by \texttt{REQUIRES\_TEST} edges. Three-hop paths are preferred over
\texttt{DIRECTLY\_INDICATES\_TEST} shortcuts because they carry the condition
context necessary for clinical validation. Triads are deduplicated on (finding,
test) pairs and sorted by clinical acuity (HIGH $\rightarrow$ MODERATE
$\rightarrow$ STANDARD) to prioritise the most time-critical pathways first.

\subsection*{Persona Descriptions}

Four personas operate sequentially on each triad. Personas A--C produce
independent assessments with no shared context at inference time; Persona D
synthesizes their outputs into a final verdict.

\paragraph{Persona A — Protocol \& Triage Enforcer.}
The only persona with external tool access. Persona A queries a local ChromaDB
vector database built from the EPFL clinical guidelines dataset, retrieving the
top-5 most relevant passages via semantic search filtered by test type. Outputs
include a structured \texttt{support\_level} (\textit{strong} $|$
\textit{moderate} $|$ \textit{weak} $|$ \textit{none}), a matched guideline
reference, whether a validated clinical decision rule is required before
ordering, and a Choosing Wisely flag.

\paragraph{Persona B — Rapid Diagnostics Researcher.}
Operates without an LLM call. Persona B queries Semantic Scholar for literature
co-evidence on the condition--test pair (with Europe PMC as fallback), returning
paper count, mean publication year (classified as \textit{current} $\geq 2020$,
\textit{aging} $\geq 2015$, \textit{outdated} $< 2015$), and keyword flags
indicating whether point-of-care alternatives or negative predictive value
evidence appear in retrieved abstracts.

\paragraph{Persona C — ED Attending.}
A pure LLM call with no tool access. Persona C evaluates clinical flow
pragmatics: over-testing risk, under-testing risk, safe discharge logic,
wrong-test detection, and whether the pathway constitutes never-event mitigation
— i.e., whether the test is mandated to catch a catastrophic, time-sensitive
diagnosis.

\paragraph{Persona D — ED Medical Director.}
Receives the full structured outputs of Personas A, B, and C and synthesizes a
final \texttt{audit\_status} (\textit{Verified} $|$ \textit{Flagged for Review}
$|$ \textit{Rejected}) and a \texttt{confidence\_score} $\in [0, 1]$. The
weighting scheme is explicit: Persona A carries the dominant weight; Persona B
adjusts confidence by up to $\pm 0.15$ based on literature recency and paper
count; Persona C's \texttt{pathway\_is\_never\_event\_mitigation} flag adds
$+0.10$, while \texttt{wrong\_test\_concern} is a hard-flag trigger regardless
of other scores.

\subsection*{Hard-Flag Triggers}

Three conditions unconditionally produce a \textit{Flagged for Review} verdict,
bypassing the scoring logic:
\begin{enumerate}
    \item Persona C raises \texttt{wrong\_test\_concern}, indicating a genuine
    clinical error in the Knowledge Graph.
    \item Persona A raises a Choosing Wisely flag.
    \item A rule-out pathway has zero guideline support and zero literature
    evidence simultaneously.
\end{enumerate}

\subsection*{Audit Metadata and Checkpointing}

Audit results are written back to the graph in-place. Each audited edge
receives \texttt{audit\_status}, \texttt{confidence\_score},
\texttt{last\_audited}, persona-level sub-scores, and \texttt{audit\_rationale}
as edge attributes, making results queryable during GraphRAG retrieval. A
complete \texttt{audit\_report.json} is also produced containing per-triad
detail and a summary of verified, flagged, and rejected counts with mean
confidence. To manage LLM call volume over a fully enumerated graph, completed
triad results are cached and skipped on re-invocation via a \texttt{--resume}
flag.

% ─────────────────────────────────────────────────────────────────────────────
\newpage
\section{Detailed Imputation Methodologies}

\subsection*{Knowledge Graph Traversal Algorithm}
\label{app:kg_matching}

Algorithm~\ref{alg:kg_imputation} formalizes the deterministic substring
matching and graph traversal procedure. The bidirectional substring check
captures clinical shorthand and abbreviations; 1-hop and 2-hop traversal
identifies both directly indicated tests and condition-mediated test
requirements.

\begin{algorithm}[htbp]
\caption{Knowledge Graph Entity Matching and Test Imputation}
\label{alg:kg_imputation}
\begin{algorithmic}[1]
\Require Patient diagnosis string $D$, Knowledge Graph $G=(V, E)$, minimum node threshold $\tau$
\Ensure Set of imputed diagnostic tests $P$

\State $T \gets \text{Split}(D, \texttt{";"})$ \Comment{Isolate terms from compound diagnosis}
\State $M \gets \emptyset$ \Comment{Matched KG nodes}

\State \textbf{Step 1: Bidirectional Substring Matching}
\For{each $t \in T$}
    \State $t \gets \text{Lowercase}(\text{Strip}(t))$
    \For{each $v \in V$ where $\text{Type}(v) \in \{\text{Condition, Symptom, Risk\_Factor, \ldots}\}$}
        \State $L \gets \text{Lowercase}(\text{Label}(v))$
        \If{$t \subseteq L$ \textbf{or} $L \subseteq t$}
            \State $M \gets M \cup \{v\}$
        \EndIf
    \EndFor
\EndFor

\State \textbf{Step 2: Graph Traversal and Test Voting}
\State $\mathit{TestCounts} \gets \text{empty map}$ \Comment{Unique nodes indicating each test}
\For{each $m \in M$}
    \State $\mathit{TestsForM} \gets \emptyset$
    \For{each outgoing edge $(m, u) \in E$}
        \If{$\text{Relation}(m,u) \in \{\texttt{REQUIRES\_TEST},\; \texttt{DIRECTLY\_INDICATES\_TEST}\}$}
            \State $\mathit{TestsForM} \gets \mathit{TestsForM} \cup \{u\}$ \Comment{1-hop}
        \ElsIf{$\text{Relation}(m,u) = \texttt{INDICATES\_CONDITION}$}
            \For{each outgoing edge $(u, w) \in E$}
                \If{$\text{Relation}(u,w) = \texttt{REQUIRES\_TEST}$}
                    \State $\mathit{TestsForM} \gets \mathit{TestsForM} \cup \{w\}$ \Comment{2-hop}
                \EndIf
            \EndFor
        \EndIf
    \EndFor
    \For{each $test \in \mathit{TestsForM}$}
        \State $\mathit{TestCounts}[test] \gets \mathit{TestCounts}[test] + 1$
    \EndFor
\EndFor

\State \textbf{Step 3: Threshold Filtering}
\State $P \gets \{test \mid \mathit{TestCounts}[test] \geq \tau\}$
\State \Return $P$
\end{algorithmic}
\end{algorithm}

\subsection*{Semantic Retrieval and Score Aggregation}
\label{app:semantic_matching}

Algorithm~\ref{alg:semantic_imputation} formalizes the GraphRAG semantic
pipeline. Patient diagnoses are embedded using the MedEmbed biomedical language
model and queried against a ChromaDB vector store. Before accepting a match, a
six-token negation window scans the original text for tokens such as \textit{no,
not, denies, without, absent, negative for, rules out}; matches under negation
are discarded. Instead of discrete voting, each matched node contributes its
cosine similarity score to the tests it indicates, and tests are filtered by an
aggregate score threshold $\theta$.

\begin{algorithm}[htbp]
\caption{GraphRAG Semantic Retrieval and Score-Based Imputation}
\label{alg:semantic_imputation}
\begin{algorithmic}[1]
\Require Patient diagnosis $D$, Knowledge Graph $G=(V,E)$, Vector Store $\mathcal{V}$, thresholds $k, s_{\min}, \theta$
\Ensure Set of imputed diagnostic tests $P$

\State $T \gets \text{Split}(D, \texttt{";"})$ \Comment{Isolate terms from compound diagnosis}
\State $M \gets \text{empty map}$ \Comment{Max cosine score per matched KG node}

\State \textbf{Step 1: Vector Search and Negation Filtering}
\For{each $t \in T$}
    \State $E_t \gets \text{Embed}(t)$
    \State $\mathit{Candidates} \gets \text{QueryVectorStore}(\mathcal{V},\, E_t,\, \text{top\_k}=k)$
    \For{each $c \in \mathit{Candidates}$}
        \State $score \gets 1.0 - c.\text{distance}$ \Comment{Distance to cosine similarity}
        \If{$score \geq s_{\min}$ \textbf{and not} $\text{IsNegated}(D,\, c.\text{label},\, \text{window}=6)$}
            \State $M[c.\text{node}] \gets \max\bigl(M.\text{get}(c.\text{node},\, 0),\; score\bigr)$
        \EndIf
    \EndFor
\EndFor

\State \textbf{Step 2: Graph Traversal and Score Aggregation}
\State $\mathit{TestScores} \gets \text{empty map}$
\For{each $(m,\, score) \in M$}
    \State $\mathit{TestsForM} \gets \emptyset$
    \For{each outgoing edge $(m, u) \in E$}
        \If{$\text{Relation}(m,u) \in \{\texttt{REQUIRES\_TEST},\; \texttt{DIRECTLY\_INDICATES\_TEST}\}$}
            \State $\mathit{TestsForM} \gets \mathit{TestsForM} \cup \{u\}$
        \ElsIf{$\text{Relation}(m,u) = \texttt{INDICATES\_CONDITION}$}
            \For{each outgoing edge $(u, w) \in E$}
                \If{$\text{Relation}(u,w) = \texttt{REQUIRES\_TEST}$}
                    \State $\mathit{TestsForM} \gets \mathit{TestsForM} \cup \{w\}$
                \EndIf
            \EndFor
        \EndIf
    \EndFor
    \For{each $test \in \mathit{TestsForM}$}
        \State $\mathit{TestScores}[test] \gets \mathit{TestScores}.\text{get}(test,\, 0) + score$
    \EndFor
\EndFor

\State \textbf{Step 3: Aggregate Thresholding}
\State $P \gets \{test \mid \mathit{TestScores}[test] \geq \theta\}$
\State \Return $P$
\end{algorithmic}
\end{algorithm}

% ─────────────────────────────────────────────────────────────────────────────
\newpage
\section{Computational Latency Benchmarks}
\label{app:latency_benchmarks}

To isolate the computational bottleneck within the GraphRAG semantic pipeline,
benchmarking separated inference into two phases: entity matching (embedding
generation and vector search) and graph traversal (score aggregation and test
imputation). Table~\ref{tab:latency_benchmarks} reports mean latency across 50
iterations for the matching phase and 25 iterations for the full end-to-end
pipeline.

The overwhelming majority of latency in the semantic pipeline occurs during
Phase 1, driven by the computational cost of the local biomedical embedding
model. The subsequent graph traversal (Phase 2) adds only ${\sim}37$\,ms
overhead, confirming that the bottleneck is the encoder rather than graph
structure.

\begin{table}[htbp]
    \centering
    \caption{Computational latency benchmarks comparing the baseline lexical
    string matching strategy against the GraphRAG semantic retrieval pipeline.}
    \label{tab:latency_benchmarks}
    \small
    \begin{tabular}{lrr}
        \toprule
        \textbf{Metric} & \textbf{String Matching} & \textbf{Semantic (GraphRAG)} \\
        \midrule
        \textbf{Cold-Start Initialization} & & \\
        \hspace{3mm} KG Pickle / Model Load & 6.16\,ms & 17{,}745.32\,ms \\
        \addlinespace
        \textbf{Per-Encounter Latency} & & \\
        \hspace{3mm} Phase 1: Entity Matching Only & 0.42\,ms & 137.92\,ms \\
        \hspace{3mm} Phase 2: Full Pipeline (w/ Traversal) & 0.55\,ms & 175.54\,ms \\
        \bottomrule
    \end{tabular}
\end{table}

% ─────────────────────────────────────────────────────────────────────────────
\newpage
\section{Diagnotix Web Application Architecture}
\label{app:diagnotix}

Diagnotix is built on a React~19, TypeScript, and Vite frontend coupled with a
Python~3.12 FastAPI backend. The system supports LLM integration through both
Nebius (Kimi-K2.5-fast) and Azure OpenAI. The application is divided into two
primary operational modes: \textit{Navigate} and \textit{Analyze}.

\subsection*{Navigate Mode}

Navigate Mode is the default visual exploration interface for the clinical
Knowledge Graph.
\begin{itemize}
    \item \textbf{Interactive Graph Canvas.} A force-directed 2D visualization
    powered by \texttt{react-force-graph-2d}.
    \item \textbf{Pathway Picker.} Filters the visualization to the 1-hop
    subgraph of a selected diagnostic test (e.g., ECG, Appendix Ultrasound). An
    ``All Pathways'' option reveals the full macroscopic graph structure.
    \item \textbf{Node Search \& Legend.} Facilitates node search by label with
    an interactive legend for toggling node-type visibility.
    \item \textbf{Dynamic Pathway Generation.} Users input a novel diagnostic
    test (e.g., CT Head), triggering a backend pipeline that grounds itself via
    RAG over Qdrant and live Semantic Scholar abstracts to generate 8--15 triage
    rules (up to 3 convergence iterations). New rules are appended to
    \texttt{guideline\_rules.json} and the Knowledge Graph is rebuilt; only the
    delta (new nodes and edges) is returned to the canvas, avoiding full page
    reloads.
\end{itemize}

\subsection*{Analyze Mode}

Analyze Mode provides a conversational interface for a selected diagnostic
pathway.
\begin{itemize}
    \item \textbf{Contextual Injection.} The selected pathway's subgraph (nodes
    and edges) is injected directly into the LLM context window.
    \item \textbf{Evidence-Grounded Generation.} Responses from Nebius
    Kimi-K2.5-fast are enriched by RAG retrieval from clinical guidelines
    (Qdrant) and live Semantic Scholar abstracts; all outputs are cited.
    \item \textbf{State Management.} Multi-turn conversation history is
    preserved independently per pathway.
\end{itemize}

\subsection*{Backend API Endpoints}

\begin{table}[htbp]
    \centering
    \caption{Diagnotix FastAPI backend endpoints.}
    \label{tab:api_endpoints}
    \small
    \resizebox{\linewidth}{!}{%
    \begin{tabular}{ll}
        \toprule
        \textbf{Endpoint} & \textbf{Purpose} \\
        \midrule
        \texttt{GET /api/tests}
          & Retrieve a list of all diagnostic test nodes. \\
        \addlinespace
        \texttt{GET /api/graph?pathway=\textless{}name\textgreater}
          & Retrieve the full graph or a specific 1-hop pathway subgraph. \\
        \addlinespace
        \texttt{POST /api/add\_test}
          & Execute the LLM pipeline to generate rules; return the node/edge diff. \\
        \addlinespace
        \texttt{POST /api/chat}
          & Initialize clinical chat with RAG and Semantic Scholar grounding. \\
        \bottomrule
    \end{tabular}}
\end{table}
